%======================================================================
\section{Diagonalization of the QCD Hamiltonian at $\order{g}$ and $\order{g^{2}}$}
\label{sec:diag}
%======================================================================

The evolution operator provides a systematic framework for diagonalizing the soft QCD
Hamiltonian of a fast-moving hadron. We use the unitary similarity transformation
\begin{equation}
H\to\OM^{\dg}H\,\OM ,
\label{eq:sim}
\end{equation}
with $\OM$ constructed perturbatively in $g$ with the coefficients fixed in
section~\ref{sec:coeff}. The role of $\OM$ is to cancel the off-diagonal matrix elements of $H$
that connect different Fock sectors, order by order in the coupling. In Paper~I it was shown that
in pure Yang--Mills theory the coherent operator \eqref{eq:coherent} alone is not sufficient: the
three-gluon interaction generates additional structures at the same order. With quarks, the
quark--gluon vertex generates three further structures at order $g$, and all of them must be
included for a complete diagonalization.

Splitting the Hamiltonian by order in the coupling,
\begin{align}
&H=H_{0}+H^{(1)}+H^{(2)},\qquad
H^{(1)}=H_{g}+H_{ggg}+H_{gq\bar q},
\nonumber\\
&H^{(2)}=H_{gggg}+H_{gg-{\rm inst}}+H_{gggg-{\rm inst}}+H_{ggq\bar q}+H_{q-{\rm inst}},
\label{eq:Hsplit}
\end{align}
and using $\OM^{(1)\dg}=-\OM^{(1)}$ and $\OM^{(2)}=iG^{(2)}+\tfrac12\big(\OM^{(1)}\big)^{2}$, one finds
\begin{align}
\OM^{\dg}H\OM\Big|_{g^{1}}&=H^{(1)}+\big[H_{0},\OM^{(1)}\big],
\label{eq:ord1}\\
\OM^{\dg}H\OM\Big|_{g^{2}}&=H^{(2)}+\tfrac12\big[H^{(1)},\OM^{(1)}\big]+\big[H_{0},iG^{(2)}\big],
\label{eq:ord2}
\end{align}
where in \eqref{eq:ord2} we used $\big[[H_{0},\OM^{(1)}],\OM^{(1)}\big]=-[H^{(1)},\OM^{(1)}]$, which
follows once \eqref{eq:ord1} vanishes. For a normal ordered structure $\mathcal O$ whose creation
operators carry total energy $E_{\rm out}$ and whose annihilation operators carry $E_{\rm in}$,
\begin{equation}
\big[H_{0},\mathcal O\big]=\big(E_{\rm out}-E_{\rm in}\big)\,\mathcal O .
\label{eq:H0O}
\end{equation}

\subsection{Order $g$}

Substituting the first order part of the ansatz into \eqref{eq:ord1} and using \eqref{eq:H0O},
\begin{align}
\OM^{\dg}H\OM\Big|_{g^{1}}=&\;H_{g}+H_{ggg}+H_{gq\bar q}
+\int_{p}\Big[A^{b}_{j}(p)\,\eps_{p}\,a^{\dg b}_{j}(p)+A^{\dg b}_{j}(p)\,\eps_{p}\,a^{b}_{j}(p)\Big]
\nonumber\\
&+\int_{p,q,r}\big(\eps_{q}+\eps_{r}-\eps_{p}\big)\Big[C(p,q,r)\,a^{\dg}(r)a^{\dg}(q)a(p)
\nonumber\\&\hspace{4cm}+C^{\dg}(r,q,p)\,a^{\dg}(r)a(q)a(p)\Big]
\nonumber\\
&+\int_{p,l,m}\big(\eps_{l}+\eps_{m}-\eps_{p}\big)\Big[\mf A_{3}(p,m,l)\,b^{\dg}(l)d^{\dg}(m)a(p)
\nonumber\\&\hspace{4cm}+\mf A^{\dg}_{3}(p,m,l)\,a^{\dg}(p)d(m)b(l)\Big]
\nonumber\\
&+\int_{p,l,l'}\big(\eps_{l}+\eps_{p}-\eps_{l'}\big)\Big[\mf F(l',l,p)\,a^{\dg}(p)b^{\dg}(l)b(l')
\nonumber\\&\hspace{4cm}+\mf F^{\dg}(l',l,p)\,b^{\dg}(l')b(l)a(p)\Big]
\nonumber\\
&+\int_{p,m,m'}\big(\eps_{m}+\eps_{p}-\eps_{m'}\big)\Big[\mf G(m',m,p)\,a^{\dg}(p)d^{\dg}(m)d(m')
\nonumber\\&\hspace{4cm}+\mf G^{\dg}(m',m,p)\,d^{\dg}(m')d(m)a(p)\Big],
\label{eq:diag1}
\end{align}
where colour, polarization and helicity labels are suppressed. The Hamiltonian $H^{(1)}$ is purely
off-diagonal --- every term of $H_{g}$, $H_{ggg}$ and $H_{gq\bar q}$ changes the Fock sector --- and the
commutator with $H_{0}$ multiplies each structure by its energy imbalance. Hence \eqref{eq:ord1} can
be solved structure by structure: writing $H_{gq\bar q}$ in the form \eqref{eq:Hqqg}, the vanishing of
\eqref{eq:diag1} requires
\begin{equation}
\mf A_{3}=\frac{V_{3}}{\eps_{p}-\eps_{l}-\eps_{m}},\qquad
\mf F=\frac{V_{F}}{\eps_{l'}-\eps_{l}-\eps_{p}},\qquad
\mf G=\frac{V_{G}}{\eps_{m'}-\eps_{m}-\eps_{p}},
\label{eq:diagsol}
\end{equation}
together with $A=-\mf a/\eps_{p}$ and $C=\mf c/(\eps_{p}-\eps_{q}-\eps_{r})$ of Paper~I, where
$\mf a$ and $\mf c$ are the vertices of $H_{g}$ and $H_{ggg}$. These are precisely the coefficients
\eqref{eq:A3}--\eqref{eq:G} obtained by matching to the LCWF: at this order the two characterizations of
$\OM$ --- it generates the dressed states and it diagonalizes $H$ --- agree term by term. Upon
substituting the explicit values, all off-diagonal contributions cancel identically and
\begin{equation}
\OM^{\dg}\big(H_{0}+H_{g}+H_{ggg}+H_{gq\bar q}\big)\OM=H_{0}+\order{g^{2}} .
\label{eq:diag1res}
\end{equation}
The transformed Hamiltonian is thus fully diagonal at order $g$. We stress that the coefficients
$\mf F$ and $\mf G$ are indispensable: without them the $a^{\dg}b^{\dg}b$ and $a^{\dg}d^{\dg}d$ pieces of
$H_{gq\bar q}$ survive in \eqref{eq:diag1} and the Hamiltonian is not diagonal at order $g$.

\subsection{Order $g^{2}$}

At order $g^{2}$ the transformed Hamiltonian is \eqref{eq:ord2}. Denote
$\mathcal K\equiv H^{(2)}+\tfrac12[H^{(1)},\OM^{(1)}]$. By \eqref{eq:H0O}, $iG^{(2)}$ can be chosen to
cancel any normal ordered piece of $\mathcal K$ with $E_{\rm out}\neq E_{\rm in}$, its coefficient being
$-\mathcal K/(E_{\rm out}-E_{\rm in})$; as shown in section~\ref{sec:lcpt}, the resulting $\OM^{(2)}$ is
exactly the LCPT coefficient $\mu_{2}$, i.e.\ the coefficients of section~\ref{sec:coeff}. What
remains is the part of $\mathcal K$ with $E_{\rm out}=E_{\rm in}$. With the $i\epsilon$ prescription the
off-diagonal elements between states of different energy cancel, and what is left for the particle
number preserving structures is multiplied by $i\epsilon/(E_{i}-E_{f}+i\epsilon)$, which is
nonvanishing only on the energy shell; since the shell is a set of measure zero, this leftover does
not contribute to matrix elements between normalizable states, exactly as in Paper~I. This applies
in particular to the $\rho$ dependent terms that rotate the colour of a soft quark in the valence
background, i.e.\ to the instantaneous $\rho\cdot J_{q}$ interaction and to the products $H_{g}\times\mf F$,
$H_{g}\times\mf G$: they transfer momentum between the soft quark and the valence charge and are removed
by $iG^{(2)}$, the remainder being the on-shell part discussed above. Substituting all the
coefficients we obtain
\begin{align}
\OM^{\dg}H\OM=&\;H_{0}-\int_{p}\eps_{p}\,A^{\dg b}_{j}(p)A^{b}_{j}(p)
\nonumber\\
&-\int_{p,q,r}\big(\eps_{q}+\eps_{r}-\eps_{p}\big)\,C^{\dg}(r,q,p)C(p,q,r)\,a^{\dg}(p)a(p)
\nonumber\\
&-\int_{p,l,m}\big(\eps_{l}+\eps_{m}-\eps_{p}\big)\,\mf A_{3}^{\dg}(p,m,l)\mf A_{3}(p,m,l)\,a^{\dg}(p)a(p)
\nonumber\\
&-\int_{l',l,p}\big(\eps_{l}+\eps_{p}-\eps_{l'}\big)\,\mf F^{\dg}(l',l,p)\mf F(l',l,p)\,b^{\dg}(l')b(l')
\nonumber\\
&-\int_{m',m,p}\big(\eps_{m}+\eps_{p}-\eps_{m'}\big)\,\mf G^{\dg}(m',m,p)\mf G(m',m,p)\,d^{\dg}(m')d(m')
\nonumber\\
&+\order{g^{3}} ,
\label{eq:diag2}
\end{align}
where in each loop term the internal momenta are fixed by the $\delta$ functions contained in the
coefficients and the colour and spin sums are implied. The first correction is the coherent
background field energy of the Weizs\"acker--Williams field. The remaining terms are the second
order light-cone energy shifts, $-\sum|{\rm vertex}|^{2}/(E_{\rm out}-E_{\rm in})$, of a single soft
gluon (from the gluon and the quark loops) and of a single soft quark and antiquark. They are
evaluated with the variables of \eqref{eq:kappa} and \eqref{eq:kappaq}: for the quark loop, using
\eqref{eq:PhiSq}, every power of $\bm\kappa$ cancels between the numerator and the energy
denominator and
\begin{equation}
\int_{l,m}\big(\eps_{l}+\eps_{m}-\eps_{p}\big)\,\mf A_{3}^{\dg}\mf A_{3}
\;\propto\;g^{2}\Nf T_{F}\int_{0}^{1}\!\frac{dz}{z(1-z)}\big[z^{2}+(1-z)^{2}\big]\int\!d^{2}\kappa ,
\end{equation}
and likewise for the gluon loop and for the quark self-energy, with the splitting functions
$P_{gg}$ and $P_{gq}$ in place of $z^{2}+(1-z)^{2}$. The transverse integrals carry no scale --- the
soft partons are massless and there is no external invariant --- and vanish in dimensional
regularization, as do the integrals $\int d^{2}\kappa/\bm\kappa^{2}$ in the normalizations
\eqref{eq:qloop} and \eqref{eq:qnorm}; the same holds for the self-energy (seagull) contributions of the
instantaneous interactions. Hence
\begin{multline}
\int_{p,q,r}\big(\eps_{q}+\eps_{r}-\eps_{p}\big)C^{\dg}C
=\int_{l,m}\big(\eps_{l}+\eps_{m}-\eps_{p}\big)\mf A_{3}^{\dg}\mf A_{3}
\\=\int_{l,p}\big(\eps_{l}+\eps_{p}-\eps_{l'}\big)\mf F^{\dg}\mf F
=\int_{m,p}\big(\eps_{m}+\eps_{p}-\eps_{m'}\big)\mf G^{\dg}\mf G=0 .
\label{eq:loops0}
\end{multline}
This is the statement that a free massless soft gluon or quark receives no mass correction, as
gauge invariance and chiral symmetry require. The remaining coherent term is
\begin{equation}
\int_{p}\eps_{p}\,A^{\dg b}_{j}(p)A^{b}_{j}(p)=\int_{p}\frac{g^{2}\,\rho^{b}(p)\,\rho^{b}(-p)}{(p^{+})^{2}},
\label{eq:WW}
\end{equation}
so that the final diagonalized Hamiltonian through $\order{g^{2}}$ is
\begin{equation}
\boxed{\;
\OM^{\dg}H\OM=H_{0}-\int_{p}\frac{g^{2}\,\rho^{a}(p)\,\rho^{a}(-p)}{(p^{+})^{2}}+\order{g^{3}},
\qquad \rho^{a}=\rho^{a}_{g}+\rho^{a}_{q}\; .}
\label{eq:diagfinal}
\end{equation}
The only surviving correction is the $\rho^{2}$ background field energy, demonstrating explicitly
that the perturbatively constructed evolution operator diagonalizes the soft QCD Hamiltonian through
$\order{g^{2}}$ within the CGC framework. Three remarks are in order.

\emph{The quark sector adds no new diagonal structure.} Every quark induced term at $\order{g^{2}}$ is
either removed by $iG^{(2)}$, because it changes the Fock sector or transfers momentum to the valence
charge, or is proportional to a scaleless integral. The form of the background field energy, and the
fact that it is the only survivor, are unchanged from pure Yang--Mills theory; the sole quark
modification of \eqref{eq:diagfinal} is that $\rho$ now contains the valence quark charge density
\eqref{eq:rhoq}. The pair creation off the vacuum is of order $g^{2}$, so its contribution to the
ground state energy is of order $g^{4}$.

\emph{Consistency with the coefficients.} The order $g$ condition reproduced $A$, $C$, $\mf A_{3}$ and
forced $\mf F$, $\mf G$; at order $g^{2}$ the construction requires $iG^{(2)}$ to cancel the off-diagonal part
of $\mathcal K$, which is the same information carried by the second order coefficients obtained by
LCPT matching. The two routes are the same calculation read in opposite directions, and the
unitarity relations of section~\ref{sec:omega} are what verifies that they agree.

\emph{Why $\Nf$ is absent.} The $\Nf$ dependent part of the running coupling originates in the quark
loop of the gluon self-energy. It does not appear in \eqref{eq:diagfinal} because this is the energy
of the soft sector in a fixed valence background with massless on-shell soft partons; the
corresponding ultraviolet pole is carried by the scaleless integral in \eqref{eq:gnorm}, where it
combines with the gluon loop into $11N_{c}-4T_{F}\Nf$. As in Paper~I, the running coupling becomes
explicit in physical observables once the $\order{g^{3}}$ coefficients are included.
